%%
%% This is file `sample-sigconf.tex',
%% generated with the docstrip utility.
%%
%% The original source files were:
%%
%% samples.dtx  (with options: `all,proceedings,bibtex,sigconf')
%%
%% IMPORTANT NOTICE:
%%
%% For the copyright see the source file.
%%
%% Any modified versions of this file must be renamed
%% with new filenames distinct from sample-sigconf.tex.
%%
%% For distribution of the original source see the terms
%% for copying and modification in the file samples.dtx.
%%
%% This generated file may be distributed as long as the
%% original source files, as listed above, are part of the
%% same distribution. (The sources need not necessarily be
%% in the same archive or directory.)
%%
%%
%% Commands for TeXCount
%TC:macro \cite [option:text,text]
%TC:macro \citep [option:text,text]
%TC:macro \citet [option:text,text]
%TC:envir table 0 1
%TC:envir table* 0 1
%TC:envir tabular [ignore] word
%TC:envir displaymath 0 word
%TC:envir math 0 word
%TC:envir comment 0 0
%%
%% The first command in your LaTeX source must be the \documentclass
%% command.
%%
%% For submission and review of your manuscript please change the
%% command to \documentclass[manuscript, screen, review]{acmart}.
%%
%% When submitting camera ready or to TAPS, please change the command
%% to \documentclass[sigconf]{acmart} or whichever template is required
%% for your publication.
%%
%%
\documentclass[sigconf,anonymous,review]{acmart}
%%
%% \BibTeX command to typeset BibTeX logo in the docs
\AtBeginDocument{%
  \providecommand\BibTeX{{%
    Bib\TeX}}}

%% Rights management information.  This information is sent to you
%% when you complete the rights form.  These commands have SAMPLE
%% values in them; it is your responsibility as an author to replace
%% the commands and values with those provided to you when you
%% complete the rights form.
\setcopyright{acmlicensed}
\copyrightyear{2018}
\acmYear{2018}
\acmDOI{XXXXXXX.XXXXXXX}
%% These commands are for a PROCEEDINGS abstract or paper.
\acmConference[Conference acronym 'XX]{Make sure to enter the correct
  conference title from your rights confirmation email}{June 03--05,
  2018}{Woodstock, NY}
%%
%%  Uncomment \acmBooktitle if the title of the proceedings is different
%%  from ``Proceedings of ...''!
%%
%%\acmBooktitle{Woodstock '18: ACM Symposium on Neural Gaze Detection,
%%  June 03--05, 2018, Woodstock, NY}
\acmISBN{978-1-4503-XXXX-X/2018/06}


%%
%% Submission ID.
%% Use this when submitting an article to a sponsored event. You'll
%% receive a unique submission ID from the organizers
%% of the event, and this ID should be used as the parameter to this command.
%%\acmSubmissionID{123-A56-BU3}

%%
%% For managing citations, it is recommended to use bibliography
%% files in BibTeX format.
%%
%% You can then either use BibTeX with the ACM-Reference-Format style,
%% or BibLaTeX with the acmnumeric or acmauthoryear sytles, that include
%% support for advanced citation of software artefact from the
%% biblatex-software package, also separately available on CTAN.
%%
%% Look at the sample-*-biblatex.tex files for templates showcasing
%% the biblatex styles.
%%

%%
%% The majority of ACM publications use numbered citations and
%% references.  The command \citestyle{authoryear} switches to the
%% "author year" style.
%%
%% If you are preparing content for an event
%% sponsored by ACM SIGGRAPH, you must use the "author year" style of
%% citations and references.
%% Uncommenting
%% the next command will enable that style.
%%\citestyle{acmauthoryear}


%%
%% end of the preamble, start of the body of the document source.
\begin{document}

%%
%% The "title" command has an optional parameter,
%% allowing the author to define a "short title" to be used in page headers.
\title{SecMutBench: Evaluating LLM-Generated Security Tests via Mutation-Based Vulnerability Detection}

%%
%% The "author" command and its associated commands are used to define
%% the authors and their affiliations.
%% Of note is the shared affiliation of the first two authors, and the
%% "authornote" and "authornotemark" commands
%% used to denote shared contribution to the research.
\author{Ben Trovato}
\authornote{Both authors contributed equally to this research.}
\email{trovato@corporation.com}
\orcid{1234-5678-9012}
\author{G.K.M. Tobin}
\authornotemark[1]
\email{webmaster@marysville-ohio.com}
\affiliation{%
  \institution{Institute for Clarity in Documentation}
  \city{Dublin}
  \state{Ohio}
  \country{USA}
}

\author{Lars Th{\o}rv{\"a}ld}
\affiliation{%
  \institution{The Th{\o}rv{\"a}ld Group}
  \city{Hekla}
  \country{Iceland}}
\email{larst@affiliation.org}

\author{Valerie B\'eranger}
\affiliation{%
  \institution{Inria Paris-Rocquencourt}
  \city{Rocquencourt}
  \country{France}
}

\author{Aparna Patel}
\affiliation{%
 \institution{Rajiv Gandhi University}
 \city{Doimukh}
 \state{Arunachal Pradesh}
 \country{India}}

\author{Huifen Chan}
\affiliation{%
  \institution{Tsinghua University}
  \city{Haidian Qu}
  \state{Beijing Shi}
  \country{China}}

\author{Charles Palmer}
\affiliation{%
  \institution{Palmer Research Laboratories}
  \city{San Antonio}
  \state{Texas}
  \country{USA}}
\email{cpalmer@prl.com}

\author{John Smith}
\affiliation{%
  \institution{The Th{\o}rv{\"a}ld Group}
  \city{Hekla}
  \country{Iceland}}
\email{jsmith@affiliation.org}

\author{Julius P. Kumquat}
\affiliation{%
  \institution{The Kumquat Consortium}
  \city{New York}
  \country{USA}}
\email{jpkumquat@consortium.net}

%%
%% By default, the full list of authors will be used in the page
%% headers. Often, this list is too long, and will overlap
%% other information printed in the page headers. This command allows
%% the author to define a more concise list
%% of authors' names for this purpose.
\renewcommand{\shortauthors}{Trovato et al.}

\newcommand{\todo}[1]{\textcolor{red}{\textit{[TODO: #1]}}}
\definecolor{editgreen}{rgb}{0,0.5,0}
\newcommand{\fix}[1]{\textcolor{editgreen}{#1}}

%%
%% The abstract is a short summary of the work to be presented in the
%% article.
\begin{abstract}
Evaluating whether Large Language Models (LLMs) can generate effective security tests remains an open challenge. Existing code security benchmarks focus on secure code \emph{generation} and rely on surface-level metrics such as pass@k, failing to measure whether generated tests actually detect vulnerabilities. We introduce SecMutBench, a mutation-based benchmark for evaluating LLM-generated security tests across \fix{30} CWE categories. SecMutBench applies \fix{25} security-specific mutation operators to \fix{339} Python programs, producing \fix{1,869} mutants that model realistic vulnerability patterns---such as converting parameterized SQL queries to string concatenation (CWE-89) or downgrading cryptographic algorithms (CWE-327). We propose the Security Mutation Score (SMS), which classifies mutant kills into semantic, \fix{functional, incidental,} and crash categories using operator-aware heuristics to distinguish genuine security awareness from superficial test failures. \fix{Our evaluation of eight LLMs alongside static analysis baselines (Bandit, Semgrep) reveals a critical finding: accounting for test validity, the best LLM achieves only 19.7\% effective SMS (SMS $\times$ Secure-Pass Rate) compared to 47.6\% for expert-written reference tests---a 2.4$\times$ gap that raw mutation scores completely obscure. Among valid test suites, traditional mutation scores further overstate security capability by an average of 2.2$\times$ compared to SMS, with functional assertions and runtime crashes---not security-aware assertions---accounting for the majority of mutant kills. Furthermore, static analysis achieves 71.4\% detection on pattern-based vulnerabilities but 0\% on logic-flaw CWEs, demonstrating complementary coverage.} SecMutBench is publicly available with an executable evaluation framework, CWE-specific mock environments, and pre-generated mutants for reproducible benchmarking.
\end{abstract}

%%
%% The code below is generated by the tool at http://dl.acm.org/ccs.cfm.
%% Please copy and paste the code instead of the example below.
%%
\begin{CCSXML}
<ccs2012>
<concept>
<concept_id>10011007</concept_id>
<concept_desc>Software and its engineering</concept_desc>
<concept_significance>500</concept_significance>
</concept>
</ccs2012>
\end{CCSXML}

\ccsdesc[500]{Software and its engineering}
%%
%% Keywords. The author(s) should pick words that accurately describe
%% the work being presented. Separate the keywords with commas.
\keywords{\fix{security benchmark, mutation testing, large language models, vulnerability detection, code security, CWE, LLM evaluation, software testing, security testing}}
%% A "teaser" image appears between the author and affiliation
%% information and the body of the document, and typically spans the
%% page.


\received{20 February 2007}
\received[revised]{12 March 2009}
\received[accepted]{5 June 2009}

%%
%% This command processes the author and affiliation and title
%% information and builds the first part of the formatted document.
\maketitle

%% ====================================================================
%% I. INTRODUCTION
%% ====================================================================
\section{Introduction}

Software vulnerabilities remain one of the most costly challenges in modern computing. The global cost of cybercrime reached \$9.22 trillion in 2024 and is projected to exceed \$10.5 trillion in 2025~\cite{cybersecventures2025}, with the average cost of a single data breach climbing to \$4.88 million globally and \$10.22 million in the United States alone~\cite{ibm2024breach}. Unpatched vulnerabilities account for approximately 20\% of all breaches, while over 29,000 critical vulnerabilities were disclosed in 2025---a 24\% increase over the prior year. Security testing is fundamental to identifying these vulnerabilities before deployment, yet manual security test creation is expensive, time-consuming, and difficult to scale.

At the same time, AI-powered coding assistants have seen explosive adoption. According to Stack Overflow's 2025 Developer Survey~\cite{stackoverflow2025survey}, 84\% of developers now use AI tools in their workflows, with 65\% using them at least weekly and 51\% of professional developers relying on them daily. Code generation has become AI's first breakout use case, growing into a \$1.9 billion ecosystem~\cite{menlo2025llm} that includes AI IDEs (Cursor~\footnote{\url{https://cursor.sh}}, Windsurf~\footnote{\url{https://codeium.com/windsurf}}), application builders (Lovable, Bolt), and enterprise coding agents (Claude Code, GitHub Copilot~\cite{pearce2022asleep}). Major technology companies report that 25\% or more of their code is now AI-generated, and Menlo Ventures estimates that code-generation-focused AI captured 42\% of the LLM developer market by mid-2025~\cite{menlo2025llm}.

This rapid adoption introduces a critical question for software security: \emph{as LLMs increasingly generate code, can they also generate effective security tests to catch vulnerabilities in that code?} The stakes are high---if AI-generated code is deployed with inadequate security testing, the same tools that accelerate development may simultaneously accelerate vulnerability introduction.

Several benchmarks have emerged to evaluate LLM security capabilities, but they focus almost exclusively on secure code \emph{generation}. CyberSecEval~\cite{bhatt2023cyberseceval} evaluates whether LLMs produce vulnerable code completions across multiple languages, using static analysis rules for automated detection. SecurityEval~\cite{siddiq2022securityeval} provides 130 prompt-based Python coding challenges mapped to 75 CWE types, while SecCodePLT~\cite{secccodeplt2024} extends this with ground-truth secure/insecure code pairs and automated evaluation pipelines. CWEval~\cite{yang2024cweval} offers 119 expert-verified tasks across 31 CWE types with CodeQL-based detection, enabling outcome-driven evaluation of both functionality and security. More recently, CFCEval~\cite{cheng2025cfceval} introduced the Element-Level Relevance Metric (ELRM) to assess code fix quality across four dimensions, addressing dataset bias through the MLVBench benchmark. Pearce et al.~\cite{pearce2022asleep} assessed the security of GitHub Copilot's code contributions, finding that approximately 40\% of generated programs contained vulnerabilities. Tony et al.~\cite{tony2023llmseceval} created natural language prompts for security evaluations across multiple CWE categories.

A critical distinction separates all these benchmarks from SecMutBench: they focus on code \emph{generation} quality---evaluating whether LLM-produced code is secure. SecMutBench instead evaluates test \emph{generation} quality---measuring whether LLM-produced tests can detect vulnerabilities. These are complementary but fundamentally different capabilities. A model that generates secure code may still be unable to articulate \emph{what makes alternative implementations vulnerable} through targeted test assertions. Furthermore, a model may avoid generating vulnerable code simply because it has learned to produce safe patterns, without possessing any deeper understanding of the vulnerability mechanism.

This gap in evaluation is compounded by a gap in \emph{metrics}. Current approaches rely on surface-level measures---pass rates, code similarity scores (e.g., CodeBLEU), and binary pass/fail against known vulnerabilities---that cannot distinguish between a test that detects a SQL injection through a targeted assertion (\texttt{assert ``parameterized'' in query\_method}) and one that merely crashes when encountering malformed input (\texttt{TypeError: unsupported operand}). This distinction is critical in practice: a crash-inducing test provides no diagnostic value to a developer, while a security-aware assertion identifies the specific vulnerability class and suggests a remediation path. Yet no existing benchmark makes this distinction.

We address these gaps with SecMutBench, a mutation-based benchmark grounded in two decades of security mutation testing research~\cite{du2000testing,shahriar2008music,mouelhi2007mutation}. Rather than evaluating whether LLM-generated code is secure, we evaluate whether LLM-generated \emph{tests} can detect security vulnerabilities. Our approach applies \fix{25} security-specific mutation operators that transform secure code into realistic vulnerable variants---such as converting parameterized SQL queries to string concatenation (CWE-89), replacing \texttt{subprocess.run([...])} with \texttt{shell=True} string commands (CWE-78), or downgrading SHA-256 to MD5 (CWE-327). Each mutant models a known vulnerability pattern mapped to a CWE category~\cite{mitre2024cwe}. LLM-generated tests are then evaluated against these mutants, and kills are classified by type---semantic, \fix{functional,} incidental, or crash---to reveal whether models demonstrate genuine security awareness or merely detect code breakage.

\fix{Our evaluation reveals two striking findings. First, only 7--47\% of LLM-generated test suites even pass on secure code; accounting for this, the best model achieves only 19.7\% effective security mutation score compared to 47.6\% for expert-written tests---a 2.4$\times$ gap that raw metrics completely obscure. Second, among valid test suites, traditional mutation scores further overstate capability by an average of 2.2$\times$, with functional assertions and runtime crashes---not security-aware assertions---accounting for the inflation. These compounding distortions mean that a model reporting 90\% mutation score may deliver only 5\% effective security detection in practice.}

Our key contributions are:

\begin{enumerate}
    \item \textbf{SecMutBench benchmark}: \fix{339 Python programs across 30 CWE categories with 25 mutation operators producing 1,869 pre-generated mutants from curated data sources (SecCodePLT, CyberSecEval, SecurityEval, CWEval, hand-crafted templates, and LLM-generated variations)}, accompanied by an executable evaluation framework with 11 CWE-specific mock environments, a \texttt{SafeOS} sandbox blocking dangerous system calls, and a multi-stage contamination prevention pipeline enabling automated, reproducible evaluation.

    \item \textbf{Security Mutation Score (SMS)}: A metric using operator-aware kill classification to categorize mutant kills into semantic (security-aware assertions), \fix{functional (behavioral-contract assertions), incidental (data-value assertions),} and crash (runtime failures), providing a more accurate measure of genuine security test quality than raw mutation scores.

    \item \textbf{Empirical evaluation}: \fix{An evaluation of eight LLMs (GPT-5.2, GPT-5-mini, GPT-oss-120B, Qwen3-Coder 30B, Qwen2.5-Coder 14B, Kimi-K2.5, GLM-5, DeepSeek-Coder-V2) alongside static analysis baselines (Bandit, Semgrep) revealing that the best LLM achieves only 19.7\% effective SMS vs.\ 47.6\% for expert-written tests, that among valid tests traditional MS overstates security capability by 2.2$\times$ on average, and that static analysis and mutation testing provide complementary coverage across vulnerability types.}
\end{enumerate}

The paper is organized as follows. Section~\ref{sec:related} introduces related work on LLM security benchmarks, mutation testing for security, and LLM-based test generation. Section~\ref{sec:framework} presents the SecMutBench framework, including dataset construction, mutation operators, kill classification, and the evaluation pipeline. Section~\ref{sec:setup} describes the experimental setup. Section~\ref{sec:results} presents results addressing our \fix{four} research questions. Section~\ref{sec:discussion} discusses implications and recommendations. Section~\ref{sec:threats} addresses threats to validity, and Section~\ref{sec:conclusion} concludes.

%% ====================================================================
%% II. RELATED WORK
%% ====================================================================
\section{Related Work}\label{sec:related}

\subsection{LLM Security Code Benchmarks}

Several benchmarks evaluate LLM capabilities in security-related code tasks. CyberSecEval~\cite{bhatt2023cyberseceval} assesses insecure code suggestions across multiple languages using static analysis for detection. SecurityEval~\cite{siddiq2022securityeval} provides 130 prompt-based Python coding challenges mapped to 75 CWE types. SecCodePLT~\cite{secccodeplt2024} provides ground-truth secure/insecure code pairs with automated evaluation. CWEval~\cite{yang2024cweval} offers 119 expert-verified tasks with CodeQL-based detection. CFCEval~\cite{cheng2025cfceval} introduces the ELRM metric to evaluate vulnerability-fixing capability across four dimensions, addressing dataset bias through the MLVBench benchmark.

Pearce et al.~\cite{pearce2022asleep} assessed the security of GitHub Copilot's contributions, and Tony et al.~\cite{tony2023llmseceval} created natural language prompts for security evaluations. A critical distinction separates all these from SecMutBench: they focus on code \emph{generation} quality, while SecMutBench evaluates test \emph{generation} quality. These are complementary but fundamentally different capabilities---a model may generate secure code without understanding why alternative implementations are vulnerable.

\subsection{Mutation Testing for Security}

Mutation testing, introduced by DeMillo et al.~\cite{demillo1978hints}, evaluates test suite adequacy by introducing small syntactic changes (mutants) into programs and measuring whether tests detect these changes. The underlying \emph{competent programmer hypothesis} assumes that programs are nearly correct, and the \emph{coupling effect hypothesis} posits that tests detecting simple faults will also detect more complex ones.

The application of mutation testing to security has a two-decade research history, surveyed broadly by Jia and Harman~\cite{jia2011analysis} and Papadakis et al.~\cite{papadakis2019mutation}. Du and Mathur~\cite{du2000testing} introduced environment fault injection for security testing, and Wimmel and J\"{u}rjens~\cite{wimmel2002specification} pioneered specification-based security mutation. Shahriar and Zulkernine developed vulnerability-specific operators: MUSIC~\cite{shahriar2008music} (9 operators for SQL injection), MUTEC~\cite{shahriar2009mutec} (11 for XSS), and MUFORMAT~\cite{shahriar2008muformat} (8 for format strings). Martin and Xie~\cite{martin2007xacml} created XACML policy mutation operators. Mouelhi et al.~\cite{mouelhi2007mutation} established that functional tests achieving high code coverage fail to achieve high security mutation scores, validating the need for security-specific adequacy measures.

Loise et al.~\cite{loise2017security} proposed the first systematic framework for security-aware mutation testing, defining 15 operators for Java targeting vulnerabilities detectable by FindBugs, including operators for SQL injection, HTTP response splitting, and XML external entities. Their work demonstrated that security-specific mutation operators can effectively evaluate the security awareness of test suites. Our mutation operators in SecMutBench extend this framework to Python and introduce additional operators targeting CWEs not covered in their work, including command injection (CWE-78), insecure deserialization (CWE-502), and SSRF (CWE-918).

Recent work includes MASC~\cite{ami2023masc} (12 operators for cryptographic API misuse), $\mu$SE~\cite{bonett2018muse} (mutation-based evaluation of Android security analyzers discovering 25 undocumented flaws), and G{\"o}rz et al.'s~\cite{gorz2023mutation} mutation analysis for fuzzer evaluation. SecMutBench builds on this foundation by adapting the security mutation paradigm to evaluate LLM-generated tests---a context not addressed in prior work.

\subsection{LLM-Based Test Generation}

LLMs have shown promise in automated test generation. Sch\"{a}fer et al.~\cite{schafer2024empirical} empirically evaluated LLMs for unit test generation, Lemieux et al.~\cite{lemieux2023codamosa} used LLMs to escape coverage plateaus, Yuan et al.~\cite{yuan2024no} evaluated ChatGPT for unit tests, and Deng et al.~\cite{deng2023large} applied LLMs as zero-shot fuzzers. Work on assertion generation and test summarization~\cite{nashid2024assertion} has advanced test understanding. Notably, Dakhel et al.~\cite{dakhel2023mutap} introduced MuTAP, which leverages mutation testing to improve LLM-generated test effectiveness by augmenting prompts with surviving mutants. While MuTAP uses mutation testing as a \emph{feedback mechanism} for iterative test improvement, SecMutBench uses it as an \emph{evaluation criterion} for measuring security awareness---a complementary but distinct objective.

Despite this growing body of work, evaluation of LLM test generation has largely focused on functional correctness (line coverage, branch coverage, pass rates) rather than security-specific properties. Concurrent work on execution-driven test benchmarking (TAM-Eval~\cite{tameval2025}, AGONETEST~\cite{agonetest2025}) and mutation-guided test generation (MUTGEN~\cite{mutgen2025}) validates the approach of using mutation signals for LLM evaluation. While these works apply traditional mutation operators to measure general test adequacy, SecMutBench targets security-specific operators mapped to CWE categories, with kill classification distinguishing security awareness from incidental detection. Similarly, $\mu$BERT~\cite{mubert2022} demonstrated that ML-generated mutants can complement hand-crafted operators; future work could explore learned security mutation operators for SecMutBench. Table~\ref{tab:comparison} summarizes how SecMutBench relates to existing evaluation platforms.

\begin{table}[t]
\caption{SecMutBench vs.\ existing evaluation platforms. \textbf{Test eval.}: evaluates test generation (not just code generation). \textbf{Dyn.}: uses dynamic execution metrics. \textbf{Kill cls.}: classifies kill types. \textbf{Sec.\ mut.}: uses security-specific mutation operators. \textbf{Contam.}: contamination prevention.}
\label{tab:comparison}
\small
\begin{tabular}{lcccccr}
\toprule
\textbf{Platform} & \rotatebox{65}{\textbf{Test eval.}} & \rotatebox{65}{\textbf{Dynamic}} & \rotatebox{65}{\textbf{Kill cls.}} & \rotatebox{65}{\textbf{Sec.\ mut.}} & \rotatebox{65}{\textbf{Contam.}} & \textbf{CWEs} \\
\midrule
CyberSecEval~\cite{bhatt2023cyberseceval} & & & & & & 8 \\
SecurityEval~\cite{siddiq2022securityeval} & & & & & & 75 \\
SecCodePLT~\cite{secccodeplt2024} & & \checkmark & & & & 27 \\
CWEval~\cite{yang2024cweval} & & \checkmark & & & & 31 \\
CFCEval~\cite{cheng2025cfceval} & & & & & & -- \\
\midrule
\textbf{SecMutBench} & \checkmark & \checkmark & \checkmark & \checkmark & \checkmark & \fix{\textbf{30}} \\
\bottomrule
\end{tabular}
\end{table}


%% ====================================================================
%% III. SECMUTBENCH FRAMEWORK
%% ====================================================================
\section{SecMutBench Framework}\label{sec:framework}

%% FIGURE PLACEHOLDER 1: Framework Overview Diagram
\begin{figure}[t]
  \centering
  \fbox{\parbox{0.93\columnwidth}{\centering\vspace{2cm}\textbf{Figure Placeholder: SecMutBench Framework Overview}\\\smallskip
  \small A pipeline diagram showing: (1) Dataset Sources $\rightarrow$ (2) Secure Code + CWE Label $\rightarrow$ (3) Mutation Operators applied $\rightarrow$ (4) Pre-generated Mutants $\rightarrow$ (5) LLM Test Generation $\rightarrow$ (6) Test Execution (secure + mutants) $\rightarrow$ (7) Kill Classification (semantic / functional / incidental / crash) $\rightarrow$ (8) SMS Computation. Include mock injection into step 6.\vspace{2cm}}}
  \caption{Overview of the SecMutBench evaluation framework. Secure programs are mutated using \fix{25} security-specific operators. LLM-generated tests are executed against both secure code and mutants, and kills are classified to compute the Security Mutation Score (SMS).}
  \label{fig:framework}
\end{figure}

\subsection{Dataset Construction}

SecMutBench consolidates samples from \fix{multiple} sources to ensure diversity in vulnerability patterns, coding styles, and complexity levels. Table~\ref{tab:dataset} summarizes the dataset composition.

\begin{table}[t]
\caption{Dataset composition by source.}
\label{tab:dataset}
\small
\begin{tabular}{lrl}
\toprule
\textbf{Source} & \textbf{Samples} & \textbf{Description} \\
\midrule
\fix{SecMutBench Originals} & \fix{75} & \fix{Hand-crafted + curated from SecCodePLT,} \\
 & & \fix{CyberSecEval, SecurityEval} \\
\fix{CWEval~\cite{yang2024cweval}} & \fix{3} & \fix{Expert-verified secure/insecure pairs} \\
\fix{SecurityEval~\cite{siddiq2022securityeval}} & \fix{3} & \fix{Prompt-based challenges} \\
\fix{LLM-Generated Variations} & \fix{258} & \fix{Structurally diverse variants of originals} \\
\midrule
\textbf{Total} & \fix{\textbf{339}} & \fix{\textbf{30 CWEs, 1,869 mutants}} \\
\bottomrule
\end{tabular}
\end{table}

%% FIGURE: Dataset Composition and CWE Distribution
\begin{figure}[t]
  \centering
  \fbox{\parbox{0.93\columnwidth}{\centering\vspace{1.2cm}\textbf{Figure: Dataset Composition}\\\smallskip
  \small Two sub-panels. (a) Pie chart: dataset composition by source (\fix{Originals 24\%, LLM Variations 76\%}). (b) Horizontal bar chart: sample count per CWE category, with bars colored by source composition. Show all \fix{30} CWEs on y-axis sorted by count descending. Annotate operator count per CWE.\vspace{1.2cm}}}
  \caption{Dataset composition by source (a) and CWE distribution with per-CWE operator coverage (b).}
  \label{fig:dataset}
\end{figure}

Each sample contains a secure implementation, an insecure variant with a known vulnerability, CWE classification, reference security tests, and pre-generated mutants. Samples span three difficulty levels: \fix{easy (136), medium (101), and hard (102)}, assigned based on code complexity and vulnerability subtlety.

\textbf{Contamination prevention.} Because \fix{several of our original sources} (SecurityEval, CyberSecEval, SecCodePLT) are publicly available and likely present in LLM training data, we apply a four-stage contamination prevention pipeline to adapted samples. \emph{(1)~Temporal filtering} excludes CVE-based samples with disclosure years before a configurable cutoff (default 2024). \emph{(2)~CWE-specific structural perturbation} transforms code patterns to prevent surface-level memorization: for example, converting \texttt{db.execute()} to cursor-based patterns for CWE-89, switching between \texttt{subprocess.run}/\texttt{Popen}/\texttt{check\_output} for CWE-78, and alternating between \texttt{os.path} and \texttt{pathlib} styles for CWE-22. \emph{(3)~AST-based identifier renaming} uses an \texttt{ast.NodeTransformer} to deterministically rename user-defined functions and variables while preserving library names, with comment stripping and string literal variation applied uniformly. \emph{(4)~N-gram contamination audit} extracts 5-gram token sequences from each sample and computes Jaccard similarity against known training corpora, flagging samples exceeding a 0.3 overlap threshold. \fix{Novel SecMutBench templates and LLM-generated variations (82\% of the dataset) are exempt from perturbation since they cannot appear in prior training data.} All perturbations are validated via \texttt{compile()} with automatic fallback to the original code if syntax is broken.

\fix{
\textbf{LLM-generated variations.} To increase dataset diversity while maintaining vulnerability validity, we generated 258 structural variations of original samples using LLMs. Each variation preserves the CWE category and vulnerability mechanism of its parent sample while introducing diverse coding patterns, variable naming, and implementation approaches. All variations undergo the same validation pipeline: secure code must compile, insecure code must contain the target vulnerability pattern, and pre-generated mutants must produce valid transformations. This augmentation strategy increases sample diversity without requiring manual expert authoring.
}

\fix{
\textbf{CWE selection.} Our 30 CWEs were selected through a principled process: we began with the OWASP Top~10~(2021) and CWE Top~25~(2024), filtered to Python-relevant vulnerabilities with active CVEs reported in the past five years, and retained only those CWEs for which we could implement at least one mutation operator with verifiable secure-to-insecure transformations. The resulting set covers injection (CWE-74, 78, 89, 94, 95), broken authentication and access control (CWE-306, 639, 732, 862, 863), sensitive data exposure (CWE-209, 319, 798), cryptographic failures (CWE-326, 327, 328, 338), input validation (CWE-20, 79, 400), XML vulnerabilities (CWE-611, 643), insecure deserialization (CWE-502), SSRF (CWE-918), open redirect (CWE-601), CSRF (CWE-352), certificate validation (CWE-295), log injection (CWE-117), mass assignment (CWE-915), and file upload (CWE-434). Each CWE has expert-verified secure/insecure pairs, hand-tuned mutation operators, and operator-specific kill classification keywords.
}
%% FIGURE: Contamination Prevention Pipeline
\begin{figure}[t]
  \centering
  \fbox{\parbox{0.93\columnwidth}{\centering\vspace{1.2cm}\textbf{Figure: Contamination Prevention Pipeline}\\\smallskip
  \small Flowchart showing: Input Samples $\rightarrow$ (1) Novel Sample Tracking (SecMutBench templates exempt; adapted samples proceed) $\rightarrow$ (2) Temporal CVE Filter (cutoff year 2024; samples with pre-cutoff CVEs removed) $\rightarrow$ (3) Perturbation Pipeline (CWE-specific structural transforms $\rightarrow$ AST identifier renaming $\rightarrow$ comment stripping $\rightarrow$ string variation $\rightarrow$ \texttt{compile()} validation with fallback) $\rightarrow$ (4) N-gram Contamination Audit (5-gram Jaccard similarity; flag if $\geq$0.3 overlap) $\rightarrow$ Decontaminated Dataset. Annotate which sub-stages apply to which CWEs.\vspace{1.2cm}}}
  \caption{Four-stage contamination prevention pipeline applied to adapted samples. Novel SecMutBench templates \fix{and LLM-generated variations} bypass perturbation.}
  \label{fig:contamination}
\end{figure}


\subsection{Security Mutation Operators}

SecMutBench defines \fix{25} mutation operators that transform secure code into vulnerable variants. Unlike traditional mutation operators that create arbitrary syntactic changes~\cite{demillo1978hints}, our operators model specific vulnerability patterns mapped to CWE categories. Table~\ref{tab:operators} presents the operator taxonomy.

Our mutation operators extend the security-aware mutation testing framework introduced by Loise et al.~\cite{loise2017security}, who defined 15 operators for Java targeting vulnerabilities detectable by FindBugs. We adapt three of their operators to Python---PSQLI (SQL injection), RHTTPO (HTTP response splitting, from their RHTTPOFC), and XXE (XML external entities, from their XMLPVXXE)---and introduce \fix{22} additional operators targeting CWEs not covered in their work, including command injection (CWE-78), insecure deserialization (CWE-502), SSRF (CWE-918), IDOR (CWE-639), \fix{weak permissions (CWE-732), log injection (CWE-117), LDAP injection (CWE-643), ReDoS (CWE-400), eval injection (CWE-94), open redirect (CWE-601), file upload bypass (CWE-434)}, CSRF removal (CWE-352), and weak randomness (CWE-338). Our operator design follows their approach: we identify secure code patterns from CWE descriptions and define transformations that introduce the corresponding vulnerability.

\begin{table}[t]
\caption{\fix{Security mutation operators (top 15 by mutant count).}}
\label{tab:operators}
\small
\begin{tabular}{llr}
\toprule
\textbf{Operator} & \textbf{Target CWEs} & \textbf{n} \\
\midrule
\fix{RVALID} & \fix{CWE-20, 79, 89} & \fix{151} \\
\fix{WEAKCRYPTO} & \fix{CWE-327, 328} & \fix{128} \\
\fix{INPUTVAL} & \fix{CWE-20} & \fix{119} \\
\fix{DESERIAL} & \fix{CWE-502, 94} & \fix{106} \\
\fix{RMAUTH} & \fix{CWE-287, 306} & \fix{105} \\
\fix{EVALINJECT} & \fix{CWE-94, 95} & \fix{98} \\
\fix{SSRF} & \fix{CWE-918} & \fix{97} \\
\fix{RENCRYPT} & \fix{CWE-319, 326} & \fix{90} \\
\fix{PATHCONCAT} & \fix{CWE-22, 73} & \fix{87} \\
\fix{MISSINGAUTH} & \fix{CWE-862, 863} & \fix{83} \\
\fix{CSRF\_REMOVE} & \fix{CWE-352} & \fix{82} \\
\fix{HARDCODE} & \fix{CWE-798} & \fix{81} \\
\fix{PSQLI} & \fix{CWE-89} & \fix{80} \\
\fix{OPENREDIRECT} & \fix{CWE-601} & \fix{63} \\
\fix{LOGINJECT} & \fix{CWE-117} & \fix{62} \\
\bottomrule
\end{tabular}
\vspace{2pt}
{\footnotesize \textit{n} = \fix{total mutants produced by operator}. \fix{Additional operators (NOCERTVALID, WEAKPERM, XXE, LDAPINJECT, WEAKKEY, WEAKRANDOM, FILEUPLOAD, INFOEXPOSE, REGEXDOS, IDOR) produce 20--60 mutants each.}}
\end{table}

Each operator implements \texttt{applies\_to(code)} for pattern matching and \texttt{mutate(code)} returning (mutated\_code, description) pairs. For example, \textsc{PSQLI} converts parameterized SQL queries (\texttt{execute(query, (param,))}) to string concatenation (\texttt{execute(f"...~\{param\}")}), analogous to the Java PSQLI operator of Loise et al.~\cite{loise2017security}, and \textsc{CMDINJECT} replaces array arguments with \texttt{shell=True} string commands.

We implement a custom mutation engine rather than extending existing Python mutation testing tools (e.g., MutPy~\cite{derezinska2014experimental}) because our security-focused operators require pattern-based transformations on source code strings rather than AST node replacements. Each operator targets a specific vulnerability class defined by its CWE mapping, ensuring that generated mutants represent realistic security flaws rather than arbitrary syntactic perturbations.

%% FIGURE PLACEHOLDER 2: Mutation Operator Example
\begin{figure}[t]
  \centering
  \fbox{\parbox{0.93\columnwidth}{\centering\vspace{1.5cm}\textbf{Figure Placeholder: Mutation Operator Example}\\\smallskip
  \small Side-by-side code diff showing a concrete mutation: (Left) Secure code with parameterized SQL query \texttt{cursor.execute(query, (user\_id,))}. (Right) Mutated code with string concatenation \texttt{cursor.execute(f"SELECT ... \{user\_id\}")}. \fix{Annotate with: operator name (PSQLI), CWE-89, and the type of kill each would produce (semantic: ``assert parameterized''; functional: ``DID NOT RAISE ValueError''; incidental: ``assert len(result) == 5''; crash: TypeError).}\vspace{1.5cm}}}
  \caption{Example of the PSQLI mutation operator transforming a parameterized SQL query (secure) into string concatenation (vulnerable, CWE-89), with illustrative kill classifications.}
  \label{fig:mutation_example}
\end{figure}

Mutants are pre-generated during dataset construction and stored in each sample's JSON record. This ensures deterministic evaluation---the same mutants are used across all evaluation runs---and eliminates variation from operator application order.

\fix{
\textbf{Mutant validity and equivalent mutants.} A standard concern in mutation testing is the presence of \emph{equivalent mutants}---mutants that are syntactically different but semantically identical to the original program, making them unkillable by any test~\cite{jia2011analysis,papadakis2019mutation}. Security mutation operators are structurally less susceptible to equivalence than traditional operators because they transform \emph{behavioral security properties}, not arbitrary syntax. For instance, \textsc{PSQLI} converts parameterized SQL queries to string concatenation, which always changes whether SQL injection is possible; \textsc{EVALINJECT} introduces \texttt{eval()}, which always enables code injection. Of 25 operators, the majority alter externally observable security behavior by construction. The remaining operators (e.g., \textsc{RVALID} for input validation removal) may produce equivalence when the removed validation is unreachable; we mitigate this by requiring each operator's \texttt{applies\_to()} method to confirm the target pattern is active in the code path, and by verifying all mutants compile and execute without immediate crash on the test harness. We report the \emph{executable mutant ratio} (mutants passing compilation and basic execution / total generated) alongside MS and SMS to quantify validity.

\fix{
\textbf{Executable mutant ratio.}
Of the 1,869 pre-generated mutants, 100\% passed both
\texttt{compile()} validation and basic execution against the test
harness without immediate crash (executable mutant ratio). This
ratio varies by operator: syntactic substitution operators
(e.g.,~\textsc{PSQLI}, \textsc{WEAKCRYPTO}, \textsc{WEAKPERM})
achieve near-100\% executability, while removal operators
(e.g.,~\textsc{RVALID}, \textsc{INPUTVAL}) occasionally produce
unreachable-path equivalents where the removed validation was
never triggered by the test harness inputs. Per-operator
executability rates and equivalent-mutant estimates based on
dynamic equivalence filtering are available in the replication
package~(\url{https://github.com/Mars-2030/secmutbench}).
}

\textbf{Mutant count.} The dataset averages 5.51~mutants per program (1,869 total across 339 programs, range 4--9 per sample), which is higher than initial versions due to expanded operator coverage. Each operator targets a \emph{specific} vulnerability pattern, and programs typically match 4--6 operators based on their CWE category and code patterns. A program with a SQL query and input validation generates PSQLI, RVALID, and INPUTVAL mutants. We deliberately avoid inflating mutant counts with syntactic variants of the same vulnerability (e.g., multiple string concatenation styles for the same query), as these would artificially boost MS without adding evaluation value. The \texttt{max\_mutants} parameter (default: 10 per sample) caps generation.
}

\subsection{Kill Classification and SMS Metric}

The central insight of SecMutBench is that not all mutant kills are equal. When a test kills a mutant (passes on secure code, fails on mutant), the kill type reveals whether the test demonstrates genuine security awareness.

Formally, let $P$ denote a secure program and $M = \{m_1, m_2, \ldots, m_k\}$ the set of mutants generated by applying security mutation operators $O = \{o_1, \ldots, o_{\fix{25}}\}$. For a test suite $T$ generated by an LLM, we define the kill predicate:
\begin{equation}
    \text{kill}(t, m) = \begin{cases} 1 & \text{if } t \text{ passes on } P \text{ and fails on } m \\ 0 & \text{otherwise} \end{cases}
\end{equation}

\noindent For each killed mutant, we classify the failure into one of \fix{five} categories:

\begin{itemize}
    \item \textbf{Semantic kill}: The test fails with an assertion error whose message contains operator-specific security indicators (e.g., ``SQL query not parameterized'' for \textsc{PSQLI} mutants). This indicates the test explicitly checks for the vulnerability.

    \item \fix{\textbf{Functional kill}: The test fails because the mutation alters \emph{control-flow behavior}---an expected exception is no longer raised, a call that should fail now succeeds, or a function returns a different type---and the test's behavioral assertion (e.g., \texttt{pytest.raises(ValueError)}, \texttt{assert isinstance(result, dict)}) catches this change. The test verifies \emph{what the code does} (raises, returns, calls) without framing the check in security terms. \textbf{Distinguishing rule:} the assertion targets a behavioral contract (exception type, return type, call occurrence), not a specific data value.}

    \item \fix{\textbf{Incidental kill}: The test fails because the mutation changes a \emph{data value} that the test happens to check---an output length differs, a numeric result changes, or a string no longer matches---and a non-security equality assertion (e.g., \texttt{assert len(result) == 64}, \texttt{assert result == expected}) catches this change. The test verifies \emph{what the code outputs}, not what it does or whether it is secure. \textbf{Distinguishing rule:} the assertion targets a concrete value or content comparison, not a behavioral contract.}

    \item \textbf{Crash kill}: The test fails due to a runtime error (ImportError, TypeError, NameError, SyntaxError, AttributeError). The mutation breaks the code in a way that causes an exception rather than being detected by a test assertion.

    \item \textbf{Other kill}: The test fails due to a non-assertion, non-crash exception not captured by the above categories. These are uncommon and excluded from the primary metrics.
\end{itemize}

\fix{The five categories are mutually exclusive by construction: semantic requires operator-specific security keywords; functional requires a behavioral-contract assertion (exception, type, call) without such keywords; incidental requires a data-value assertion without such keywords; crash requires no assertion at all (runtime exception); and other captures residual cases. A kill is assigned to the first matching category in this priority order: crash $\rightarrow$ semantic $\rightarrow$ functional $\rightarrow$ incidental $\rightarrow$ other.}

\noindent Let $K_s$, \fix{$K_f$,} $K_i$, $K_c$, and $K_o$ denote the sets of semantic, \fix{functional,} incidental, crash, and other kills respectively, where $K_s \cup K_i \cup \fix{K_f \cup} K_c \cup K_o \subseteq M$ and the sets are mutually disjoint. We define the \textbf{Security Mutation Score} as:
\begin{equation}
    \text{SMS} = \frac{|K_s|}{|M|}
\end{equation}

\noindent In contrast, the traditional Mutation Score (MS) counts all kills:
\begin{equation}
    \text{MS} = \frac{|K_s| + |K_i| + \fix{|K_f| +} |K_c| + |K_o|}{|M|}
\end{equation}

\noindent The \textbf{inflation ratio} $\rho$ captures the overstatement:
\begin{equation}
    \rho = \frac{\text{MS}}{\text{SMS}}
\end{equation}

\noindent When $\rho \gg 1$, the gap is attributable to crash, functional, incidental, and other kills that do not represent genuine security detection.

\fix{Because SMS is computed only over samples where the generated test suite passes on secure code, it does not penalize models for producing invalid tests. To capture practical end-to-end effectiveness, we define the \textbf{Effective Security Mutation Score}:
\begin{equation}
    \text{EffSMS} = \text{SMS} \times \text{SPR}
\end{equation}
\noindent where SPR (Secure-Pass Rate) is the proportion of samples for which the generated test suite passes on the secure code. EffSMS is the primary comparability metric: it reflects how much security-aware detection a model delivers \emph{across the entire benchmark}, accounting for both test validity and security awareness. We additionally report SMS and SPR separately for diagnostic decomposition.}

\textbf{Vulnerability Detection (VD).} As a ground-truth sanity check, we compute a binary Vulnerability Detection metric that verifies whether generated tests can distinguish secure from insecure code:
\begin{equation}
    \text{VD}(p) = \begin{cases} 1 & \text{if } T_p \text{ passes on } P_p^{\text{sec}} \text{ and fails on } P_p^{\text{insec}} \\ 0 & \text{otherwise} \end{cases}
\end{equation}

\noindent and report the aggregate $\text{VD} = |\{p : \text{VD}(p) = 1\}| / |P|$. Note that VD is \emph{not} a mutation-based metric---it tests against the original insecure code, not mutants. Mutant-based vulnerability detection is captured by MS and SMS.

\textbf{Operator-aware classification.} A key refinement in SecMutBench is that kill classification uses operator-specific keyword lists rather than a flat global list. For example, ``path'' triggers semantic classification only for \textsc{PATHCONCAT} mutants, not for \textsc{PSQLI} mutants. Each of the \fix{25} operators has a tailored set of security-relevant terms (e.g., \textsc{WEAKCRYPTO}: \{md5, sha1, weak, bcrypt, salt\}; \textsc{DESERIAL}: \{pickle, yaml, safe\_load, SafeLoader\}). This prevents false inflation of SMS scores from coincidental keyword matches. \fix{Complete operator specifications, keyword lists, and LLM-as-Judge evaluation prompts are available in the replication package at \url{https://github.com/Mars-2030/secmutbench}.}


\subsection{Evaluation Pipeline}

The evaluation pipeline executes LLM-generated tests in isolated subprocesses using pytest. Each test execution creates a temporary directory containing the code under test, generated tests, and a conftest module that injects CWE-specific mock environments.

Mock injection serves dual purposes: \emph{safety} (preventing real system calls, file operations, and network access) and \emph{observability} (mock objects track whether security-relevant operations were performed). Eleven mock environments cover the major CWE categories: MockDatabase (CWE-89), MockSubprocess (CWE-78), MockFileSystem (CWE-22), MockHTTPClient (CWE-918), MockCrypto (CWE-327), MockPickle and MockYAML (CWE-502), MockEnvironment (CWE-798), MockXMLParser (CWE-611), MockAuthenticator (CWE-287), and MockEval (CWE-94). An additional \texttt{SafeOS} wrapper blocks dangerous operations (\texttt{os.system}, \texttt{os.popen}, \texttt{os.exec*}) to prevent LLM-generated tests from executing arbitrary system commands. Mocks are injected via both Python builtins namespace and \texttt{sys.modules} patching, ensuring that both implicit name resolution and explicit imports resolve to mock objects.

For each sample, the pipeline: (1) runs generated tests against secure code (tests should pass), (2) runs tests against each pre-generated mutant (security-aware tests should fail), (3) classifies each kill using operator-aware heuristics, and (4) computes SMS and supporting metrics. Test results are captured via a custom pytest \texttt{ResultCollector} plugin producing structured JSON.

%% FIGURE: Execution Sandbox Architecture
\begin{figure}[t]
  \centering
  \fbox{\parbox{0.93\columnwidth}{\centering\vspace{1.2cm}\textbf{Figure: Execution Sandbox Architecture}\\\smallskip
  \small Diagram showing the test execution environment. Outer box: Isolated subprocess (temp directory). Inner layers: (1) \texttt{sys.modules} patches (subprocess, requests, hashlib, pickle, yaml, os, jwt, bcrypt, flask, mysql $\rightarrow$ mock objects). (2) \texttt{builtins} namespace injection (db, fs, env, http, xml, auth, crypto, eval $\rightarrow$ mock instances). (3) \texttt{SafeOS} wrapper blocking \texttt{os.system}, \texttt{os.popen}, \texttt{os.exec*}. Center: pytest with \texttt{ResultCollector} plugin $\rightarrow$ structured JSON. Show arrows from LLM-generated test importing ``subprocess'' being intercepted by \texttt{sys.modules} to MockSubprocess.\vspace{1.2cm}}}
  \caption{Sandboxed test execution architecture showing dual mock injection (builtins + \texttt{sys.modules}), the \texttt{SafeOS} safety layer, and structured result capture.}
  \label{fig:sandbox}
\end{figure}


%% ====================================================================
%% IV. EXPERIMENTAL SETUP
%% ====================================================================
\section{Experimental Setup}\label{sec:setup}

\subsection{Research Questions}

\begin{description}
    \item[RQ1:] \textit{How effectively do LLMs generate security-aware tests?} We compare overall mutation scores (MS) against Security Mutation Scores (SMS) to quantify the inflation introduced by \fix{functional, incidental, and} crash kills.

    \item[RQ2:] \textit{What types of kills do LLM-generated tests produce?} We analyze the kill classification breakdown to understand why traditional metrics overstate LLM security testing capability.

    \item[RQ3:] \textit{How does security test effectiveness vary across vulnerability types?} We examine per-CWE and per-operator effectiveness to identify systematic blind spots, including comparison with static analysis baselines.

\textcolor{red}{
    \item[RQ4:] \textit{How sensitive is SMS to prompt design and how do LLMs compare to human-written tests?} We ablate prompt detail levels (no hint, CWE ID only, full context) and compare LLM-generated tests against expert-written reference tests to establish upper bounds. }

\end{description}


\subsection{Subject Models}

\fix{We evaluate eight LLMs spanning open-weight and proprietary models, plus two static analysis baselines. Table~\ref{tab:models} summarizes the models. All LLMs are evaluated in a zero-shot setting with temperature 0 (greedy decoding) and maximum 2048 output tokens.}

\begin{table}[t]
\caption{Subject models and baselines.}
\label{tab:models}
\small
\fix{
\begin{tabular}{lllll}
\toprule
\textbf{Model} & \textbf{Version / ID} & \textbf{Size} & \textbf{Provider} & \textbf{Access} \\
\midrule
GPT-oss-120B & \texttt{gpt-oss-120b} & 120B & OpenAI & API \\
GPT-5.2 & \texttt{gpt-5.2-2025-12-11} & -- & OpenAI & API \\
GPT-5-mini & \texttt{gpt-5-mini-2025-08-07} & -- & OpenAI & API \\
Qwen3-Coder & \texttt{qwen3-coder:30b} & 30B & Alibaba & Ollama \\
Qwen2.5-Coder & \texttt{qwen2.5-coder:14b-instruct} & 14B & Alibaba & Ollama \\
Kimi-K2.5 & \texttt{kimi-k2.5:cloud} & -- & Moonshot & Ollama \\
GLM-5 & \texttt{glm-5:cloud} & -- & Zhipu AI & Ollama \\
DeepSeek-Coder-V2 & \texttt{deepseek-coder-v2:latest} & 236B MoE & DeepSeek & Ollama \\
\midrule
Bandit & 1.9.2 & -- & Static & CLI \\
Semgrep & -- & -- & Static & CLI \\
\bottomrule
\end{tabular}
}
\end{table}

\fix{Open-weight models are served locally via Ollama providing an OpenAI-compatible API. Proprietary models are accessed via their respective cloud APIs. All kill classifications are validated using an LLM-as-Judge (GPT-5.4) to reclassify heuristic labels.}


\subsection{Prompting Strategy}

Each model receives a standardized zero-shot prompt containing the secure code implementation, the CWE category identifier and name, and the function entry point. The prompt instructs the model to generate pytest-compatible security tests that would detect if the code were vulnerable to the specified CWE. No few-shot examples are provided to establish a zero-shot baseline.

\fix{
Figure~\ref{fig:prompt} shows the full-context prompt template (abbreviated). The prompt provides: (1)~the secure code, (2)~the CWE identifier and name, (3)~detailed mock environment documentation with detection patterns for each CWE-specific mock (e.g., \texttt{db.last\_params is not None} for SQL injection, \texttt{subprocess.last\_shell == False} for command injection), (4)~an assertion logic table contrasting correct vs.\ incorrect assertion patterns, and (5)~requirements for pytest-compatible output. The no-hint variant omits all security context (``Write tests for the following Python function''), while the CWE-ID-only variant includes only the CWE identifier (``Write tests to detect \{CWE\} vulnerabilities'') without mock documentation or assertion guidance.

\begin{figure}[t]
  \centering
  \fbox{\parbox{0.93\columnwidth}{\small
  \texttt{Generate security tests for this Python function to detect \{cwe\_name\} (\{cwe\}) vulnerabilities.}\\[4pt]
  \texttt{```python}\\
  \texttt{\{code\}}\\
  \texttt{```}\\[4pt]
  \texttt{**Test Environment (already available, do NOT import):**}\\
  \textit{[Mock documentation for 11 CWE-specific mocks: MockDatabase, MockSubprocess, MockFileSystem, MockCrypto, MockPickle/YAML, MockHTTPClient, MockXMLParser, MockEnvironment, MockAuthenticator, MockEval, MockRandom]}\\[4pt]
  \texttt{**Requirements:**}\\
  \texttt{1. Write pytest-style functions starting with test\_}\\
  \texttt{2. Tests should PASS on secure, FAIL on vulnerable}\\
  \texttt{3. Test at least 3 attack vectors for \{cwe\}}\\
  \texttt{4. Use assertions to verify security properties}\\[4pt]
  \texttt{**CRITICAL: Assertion Logic for Security Tests**}\\
  \textit{[Table: XSS $\rightarrow$ assert ``<script>'' NOT in result; SQLi $\rightarrow$ assert db.last\_params is not None; CmdInj $\rightarrow$ assert subprocess.last\_shell == False; PathTrav $\rightarrow$ assert fs.last\_path.startswith(base\_dir)]}\\[4pt]
  \texttt{Output ONLY valid Python test functions.}
  }}
  \caption{Abbreviated full-context prompt template. The complete prompt includes detailed mock API documentation for all 11 mock environments.}
  \label{fig:prompt}
\end{figure}
}

\textcolor{red}{
\textbf{Prompt ablation.} To quantify how much LLM security testing depends on being told \emph{what} to look for, we evaluate three prompt variants: (a)~\emph{No hint}: ``Write security tests for this function'' (no CWE information), (b)~\emph{CWE ID only}: ``Write tests to detect CWE-89 vulnerabilities in this function,'' and (c)~\emph{Full context}: the complete prompt with CWE name, ID, and description (our default). This ablation mirrors SecCodePLT\'s~\cite{secccodeplt2024} security policy reminder experiments and reveals the extent to which SMS improvements require explicit vulnerability guidance.
}


\subsection{Baselines}

We include three baselines: (1) \textbf{Bandit}~\cite{bandit2024}, a widely-used Python static analysis tool, run on insecure code to measure detection rate per CWE; (2) \textbf{Semgrep}~\cite{semgrep2024} with the \texttt{p/security-audit} ruleset, providing a second static analysis comparison with broader rule coverage; and (3) \textbf{Reference tests}, the expert-written security tests included in each SecMutBench sample, establishing an upper-bound baseline for what hand-crafted security tests achieve on our mutants.

\fix{
Bandit detects 77 of 339 samples (22.7\% overall detection rate). On the 11 static-analyzable CWEs (147 samples), Bandit achieves 71.4\% detection; on the remaining 19 dynamic-only CWEs (192 samples), Bandit achieves 0\% detection. Semgrep with \texttt{p/python} rules detects 37 samples (10.9\%). Reference tests achieve MS of 57.7\% and SMS of 47.6\% (889 semantic kills out of 1,869 mutants), establishing the expert upper bound.
}


\subsection{Execution Environment}

\fix{All experiments run on macOS (Darwin 24.6.0) with Python~3.11.5 (required for MutPy compatibility). Each test suite execution is sandboxed in an isolated subprocess with a 5-second timeout. Open-weight models are served locally via Ollama; proprietary models are accessed via cloud APIs. Each model evaluates all 339 samples across three prompt variants (full context, CWE-ID only, no hint) in a single run with temperature 0 (greedy decoding), producing deterministic outputs. Samples are evaluated in stratified random order (proportional CWE sampling) with seed 42 for reproducibility.}

\textcolor{red}{
\textbf{Test validity.} Before computing mutation metrics, we report the \emph{test validity rate}: the proportion of LLM-generated test suites that pass on the secure code (i.e., are syntactically correct and functionally sound). Tests failing on secure code are excluded from mutation evaluation, as they cannot distinguish secure from insecure behavior. We decompose invalid tests into syntax errors, import failures, and runtime errors to diagnose generation quality.

\fix{Secure-pass rates (proportion of tests passing on secure code) are: Qwen3-Coder~30B (46.6\%), Kimi-K2.5 (42.8\%), GPT-oss-120B (36.0\%), GLM-5 (33.0\%), GPT-5-mini (29.2\%), Qwen2.5-Coder~14B (16.8\%), GPT-5.2 (13.3\%), and DeepSeek-Coder-V2 (7.1\%). The dominant failure modes are runtime errors from incorrect mock API usage (e.g., calling nonexistent mock methods), import errors from attempting to import real libraries intercepted by \texttt{sys.modules}, and assertion errors on secure code indicating inverted assertion logic. Larger open-weight models (Qwen3, Kimi) tend to produce more valid tests than smaller or API models, likely due to better adherence to the mock environment documentation.}

\fix{\textbf{Determinism.} All models use temperature 0 (greedy decoding), producing deterministic outputs for a given input. We evaluate each model once per prompt variant (3 variants $\times$ 339 samples = 1,017 evaluations per model). Because greedy decoding eliminates sampling variance, repeated runs would yield identical results; we therefore report single-run metrics without confidence intervals. Variance across prompt variants provides a sensitivity measure: the mean SMS standard deviation across the three prompt variants is 11.2 percentage points (range: 1.9 pp for DeepSeek-Coder-V2 to 23.9 pp for Qwen2.5-Coder), indicating that prompt design is the dominant source of variation rather than stochastic sampling.}
}

%% ====================================================================
%% V. RESULTS
%% ====================================================================
\section{Results}\label{sec:results}

\subsection{RQ1: Overall Effectiveness}

\fix{Table~\ref{tab:rq1} presents the overall results. Effective SMS (EffSMS = SMS $\times$ SPR) is the primary comparability metric, capturing end-to-end security testing effectiveness across the entire benchmark. SMS and SPR are reported separately for diagnostic decomposition.}

\begin{table}[t]
\caption{\fix{Overall results (OpenAI-judged, full-context prompt). EffSMS = SMS $\times$ SPR is the primary metric.}}
\label{tab:rq1}
\small
\fix{
\begin{tabular}{lrrrrr}
\toprule
\textbf{Model} & \textbf{MS} & \textbf{SPR} & \textbf{EffSMS} & \textbf{SMS} & \textbf{MS/SMS} \\
\midrule
Qwen3-Coder 30B & 75.2 & 46.6 & \textbf{19.7} & 42.3 & 1.78 \\
Kimi-K2.5 & 84.2 & 42.8 & 19.6 & 45.7 & 1.84 \\
GPT-oss-120B & 90.0 & 36.0 & 15.6 & 43.2 & 2.08 \\
GPT-5-mini & 84.2 & 29.2 & 11.6 & 39.8 & 2.12 \\
Qwen2.5-Coder 14B & 83.6 & 16.8 & 9.8 & 58.6 & 1.43 \\
GLM-5 & 40.0 & 33.0 & 9.0 & 27.4 & 1.46 \\
GPT-5.2 & 88.7 & 13.3 & 5.1 & 38.3 & 2.32 \\
DeepSeek-Coder-V2 & 51.7 & 7.1 & 0.8 & 10.9 & 4.74 \\
\midrule
Reference Tests & 57.7 & 100.0 & \textbf{47.6} & 47.6 & 1.21 \\
Bandit & 22.7 & -- & -- & -- & -- \\
\bottomrule
\end{tabular}
}
\vspace{2pt}
{\footnotesize \fix{Models sorted by EffSMS (descending). \textbf{MS}~=~Mutation Score (all kills). \textbf{SPR}~=~Secure-Pass Rate (\% of tests passing on secure code). \textbf{EffSMS}~=~Effective SMS (SMS $\times$ SPR). \textbf{SMS}~=~Security Mutation Score (semantic kills only, among valid tests). All values are percentages.}}
\end{table}

\fix{
The headline result is the gap between expert-written and LLM-generated tests. Reference tests achieve 47.6\% EffSMS; the best LLM (Qwen3-Coder~30B) achieves 19.7\%---a 2.4$\times$ gap that raw MS completely obscures (reference MS of 57.7\% is \emph{lower} than six of eight LLMs). The ranking by EffSMS differs strikingly from the ranking by MS: GPT-5.2, which ranks second on MS (88.7\%), drops to seventh on EffSMS (5.1\%) because only 13.3\% of its tests are valid. Conversely, Qwen3-Coder ranks sixth on MS but first on EffSMS, because its 46.6\% SPR compensates for moderate SMS.

Among valid test suites, the inflation ratio (MS/SMS) ranges from 1.43 (Qwen2.5-Coder) to 4.74 (DeepSeek-Coder-V2), with a mean of 2.2$\times$. Qwen2.5-Coder achieves the highest SMS (58.6\%)---exceeding reference tests (47.6\%)---but its 16.8\% SPR yields only 9.8\% EffSMS, demonstrating that raw SMS without test validity is misleading.
}

\smallskip\noindent\fbox{\parbox{0.95\columnwidth}{\fix{\textbf{Finding 1a (generation quality):} Only 7--47\% of LLM-generated test suites pass on secure code. Accounting for this, the best EffSMS is 19.7\% (Qwen3-Coder), compared to 47.6\% for expert-written reference tests---a 2.4$\times$ gap that raw mutation scores completely obscure.}}}\smallskip

\smallskip\noindent\fbox{\parbox{0.95\columnwidth}{\fix{\textbf{Finding 1b (when tests are valid):} Among valid test suites, traditional MS overstates security testing capability by 2.2$\times$ on average (range 1.4--4.7$\times$). The best model when tests are valid is Qwen2.5-Coder (58.6\% SMS), but its 16.8\% SPR makes its practical EffSMS only 9.8\%.}}}\smallskip


\subsection{RQ2: Kill Classification Analysis}

\fix{Table~\ref{tab:kills} presents the kill classification breakdown for each model's full-context prompt variant. Kill counts reflect only the mutants from samples where the model's tests passed on secure code.}

%% FIGURE PLACEHOLDER 3: Kill Classification Breakdown
\begin{figure}[t]
  \centering
  \fbox{\parbox{0.93\columnwidth}{\centering\vspace{2cm}\textbf{Figure Placeholder: Kill Classification Breakdown}\\\smallskip
  \small Stacked bar chart with one bar per model (+ Bandit baseline). Each bar divided into \fix{five} colored segments: semantic kills (green), \fix{functional kills (blue),} incidental kills (yellow), crash kills (red)\fix{, other (gray)}. X-axis: models. Y-axis: proportion of total mutants. Add a horizontal dashed line showing average SMS across models to highlight the gap between MS (full bar height) and SMS (green segment only).\vspace{2cm}}}
  \caption{Kill classification breakdown across evaluated models. The gap between total bar height (MS) and the semantic segment (SMS) illustrates the inflation ratio.}
  \label{fig:kills}
\end{figure}

\fix{
\begin{table}[t]
\caption{Kill classification breakdown (full-context prompt, counts).}
\label{tab:kills}
\small
\begin{tabular}{lrrrrrr}
\toprule
\textbf{Model} & \textbf{Mut.} & \textbf{Sem.} & \textbf{Func.} & \textbf{Incid.} & \textbf{Crash} & \textbf{Other} \\
\midrule
GPT-oss-120B & 672 & 290 & 194 & 22 & 54 & 45 \\
GPT-5.2 & 266 & 102 & 87 & 15 & 12 & 20 \\
GPT-5-mini & 538 & 214 & 155 & 18 & 22 & 44 \\
Kimi-K2.5 & 795 & 363 & 171 & 23 & 66 & 46 \\
Qwen2.5 14B & 324 & 190 & 37 & 26 & 14 & 4 \\
Qwen3 30B & 839 & 355 & 100 & 53 & 69 & 54 \\
DeepSeek-V2 & 147 & 16 & 37 & 0 & 19 & 4 \\
GLM-5 & 610 & 167 & 41 & 2 & 26 & 8 \\
\midrule
Reference & 1869 & 889 & 0 & 60 & 73 & 56 \\
\bottomrule
\end{tabular}
\vspace{2pt}
{\footnotesize \textbf{Mut.}~=~mutants evaluated (from valid-test samples). \textbf{Sem.}~=~semantic kills. \textbf{Func.}~=~functional kills. \textbf{Incid.}~=~incidental. Reference tests evaluate all 1,869 mutants (100\% SPR).}
\end{table}
}

\fix{
A key finding is the emergence of \textbf{functional kills} as a dominant non-semantic category for LLM-generated tests. \fix{Functional kills occur when tests detect behavioral changes caused by mutations through behavioral-contract assertions (e.g., expected exceptions not raised, return type mismatches) without any security reasoning.} For GPT-oss-120B, functional kills (194) account for 32.1\% of all kills---more than crash kills (54, 8.9\%). This contrasts sharply with reference tests, which produce zero functional kills, indicating that expert-written security tests are precisely targeted at vulnerability mechanisms.

Crash kills remain significant but are not the dominant non-semantic category for most models. Qwen3-Coder produces the most crash kills (69) and the most incidental kills (53), suggesting its tests interact broadly with the code but without focused security assertions. Qwen2.5-Coder stands out with the lowest non-semantic kill proportion: only 81 non-semantic kills out of 271 total (29.9\%), indicating more targeted security assertions.
}

\textcolor{red}{

\textbf{Concrete examples.} To illustrate the kill taxonomy, we present three representative kills from our evaluation:



\begin{itemize}

    \item \fix{\textbf{Semantic kill (WEAKCRYPTO, CWE-327):} A Kimi-K2.5 test asserts \texttt{not hashlib.weak\_algorithm\_used} after calling a hashing function. The WEAKCRYPTO mutant substitutes \texttt{sha256} with \texttt{md5}, causing the mock to set \texttt{weak\_algorithm\_used = True}. The assertion fails with a message referencing the weak algorithm---classified as semantic via mock observability.}

    \item \fix{\textbf{Functional kill (MISSINGAUTH, CWE-862):} A test expects \texttt{PermissionError} when a non-admin user calls a privileged function. The MISSINGAUTH mutant bypasses the \texttt{.is\_admin} check, allowing the call to succeed. The assertion \texttt{DID NOT RAISE <class 'PermissionError'>} catches the mutation, but without security-specific terminology in the error message.}

    \item \fix{\textbf{Crash kill (XXE, CWE-611):} A test calls an XML parsing function. The XXE mutant replaces \texttt{defusedxml} with \texttt{minidom}, causing \texttt{AttributeError: 'Document' object has no attribute 'get'} due to API mismatch. No assertion is evaluated---the code crashes before reaching test logic.}

    \item \fix{\textbf{Incidental kill (REGEXDOS, CWE-400):} A test asserts a file read size limit: \texttt{assert 104857600 == (1048576 + 1)}. The REGEXDOS mutant removes the size limit, changing the return value. The assertion is a concrete data-value comparison that happens to catch the mutation, not a security-aware DoS detection.}

\end{itemize}



\fix{\textbf{Classification audit.} Rather than relying solely on heuristic keyword matching, we validate all kill classifications using an LLM-as-Judge (GPT-5.4). The judge receives each kill's error message, operator, CWE, and test code, and independently classifies it into semantic, functional, incidental, crash, or other. All results reported in this paper use the judge-validated classifications. We additionally performed a human validation study on a stratified sample of 120 kills to assess heuristic--annotator agreement, yielding Cohen's $\kappa = 0.502$ (moderate agreement).

Table~\ref{tab:ira} shows per-operator agreement, revealing a
systematic pattern rather than random noise: \emph{syntactic}
operators (PSQLI, WEAKPERM, LOGINJECT, DESERIAL) achieve 89--100\%
agreement, while \emph{behavioral} operators (IDOR, OPENREDIRECT,
MISSINGAUTH, SSRF) achieve 0--17\% agreement. This split directly
reflects whether the kill's security intent is visible in the error
message (syntactic operators produce assertion failures with
security-relevant terms) or only in the test logic (behavioral
operators produce \texttt{DID NOT RAISE} failures with no security-specific
content). The disagreement is therefore \emph{structurally
predictable} from operator type, not from annotator subjectivity.

\begin{table}[t]
\caption{Per-operator annotator--heuristic agreement on kill
classification. Syntactic operators (top) show near-perfect
agreement; behavioral operators (bottom) show systematic divergence
driven by the \texttt{DID NOT RAISE} classification gap.}
\label{tab:ira}
\small
\setlength{\tabcolsep}{4pt}
\begin{tabular}{llrrr}
\toprule
\textbf{Operator} & \textbf{Type} &
\textbf{n} &
\textbf{Annot.\ sem.} &
\textbf{Agreement} \\
\midrule
PSQLI        & Syntactic &  7 & 100\% & 100\% \\
WEAKPERM     & Syntactic &  9 & 100\% & 100\% \\
LOGINJECT    & Syntactic &  3 & 100\% & 100\% \\
DESERIAL     & Syntactic & 18 &  17\% &  89\% \\
WEAKRANDOM   & Syntactic &  4 &  75\% &  75\% \\
\midrule
RVALID       & Behavioral & 11 &  27\% &  64\% \\
PATHCONCAT   & Behavioral &  3 &  33\% &  67\% \\
SSRF         & Behavioral & 13 &  46\% &  54\% \\
EVALINJECT   & Behavioral &  8 &  13\% &  50\% \\
XXE          & Behavioral &  4 &  50\% &  50\% \\
INPUTVAL     & Behavioral & 15 &  53\% &  47\% \\
MISSINGAUTH  & Behavioral &  6 & 100\% &  17\% \\
OPENREDIRECT & Behavioral &  4 & 100\% &   0\% \\
IDOR         & Behavioral &  7 & 100\% &   0\% \\
\bottomrule
\end{tabular}
\vspace{3pt}
{\footnotesize \textbf{Annot.\ sem.}: proportion of kills the
annotator labeled semantic. \textbf{Agreement}: proportion where
annotator and heuristic agree. For behavioral operators where
annotator labels semantic but heuristic labels functional (e.g.,
SSRF, IDOR, OPENREDIRECT), the heuristic is conservative: it misses
semantic kills whose security intent is expressed in the test logic
rather than the error message. The LLM-as-Judge reclassification
corrects this conservative bias for the reported results.}
\end{table}
}
}

\smallskip\noindent\fbox{\parbox{0.95\columnwidth}{\fix{\textbf{Finding 2:} Functional kills---not crash kills---are the dominant non-semantic category for LLM-generated tests, accounting for 15--36\% of all kills. Reference tests produce zero functional kills. This reveals that LLMs often detect mutations through behavioral side-effects rather than security-targeted assertions.}}}\smallskip


\subsection{RQ3: Vulnerability-Specific Analysis}

\fix{Table~\ref{tab:rq3} presents per-CWE effectiveness aggregated across models, alongside Bandit detection rates. CWEs are grouped by whether they are amenable to static analysis (pattern-based) or require dynamic behavioral testing (logic-flaw).}

\begin{table}[t]
\caption{Per-CWE effectiveness \fix{(selected CWEs, aggregated across models)}.}
\label{tab:rq3}
\small
\begin{tabular}{llrrr}
\toprule
\textbf{CWE} & \textbf{Category} & \textbf{n} & \fix{\textbf{Mut.}} & \textbf{Bandit} \\
\midrule
\fix{CWE-89} & \fix{SQL Injection} & \fix{16} & \fix{80} & \fix{71.4\%} \\
\fix{CWE-22} & \fix{Path Traversal} & \fix{15} & \fix{87} & \fix{71.4\%} \\
\fix{CWE-94} & \fix{Code Injection} & \fix{15} & \fix{102} & \fix{71.4\%} \\
\fix{CWE-327} & \fix{Weak Crypto} & \fix{12} & \fix{71} & \fix{71.4\%} \\
\fix{CWE-798} & \fix{Hardcoded Creds} & \fix{13} & \fix{81} & \fix{71.4\%} \\
\midrule
\fix{CWE-306} & \fix{Missing Auth} & \fix{16} & \fix{97} & \fix{0\%} \\
\fix{CWE-352} & \fix{CSRF} & \fix{14} & \fix{82} & \fix{0\%} \\
\fix{CWE-862} & \fix{Missing Authz} & \fix{10} & \fix{62} & \fix{0\%} \\
\fix{CWE-639} & \fix{IDOR} & \fix{5} & \fix{20} & \fix{0\%} \\
\fix{CWE-918} & \fix{SSRF} & \fix{15} & \fix{97} & \fix{0\%} \\
\bottomrule
\end{tabular}
\vspace{2pt}
{\footnotesize \fix{Top group: static-analyzable CWEs (Bandit achieves 71.4\% detection). Bottom group: logic-flaw CWEs (Bandit achieves 0\%). \textbf{n}~=~samples, \textbf{Mut.}~=~total mutants. \fix{Full 30-CWE breakdown available in the replication package~(\url{https://github.com/Mars-2030/secmutbench}).}}}
\end{table}

\fix{The static-analyzable vs.\ logic-flaw divide is stark: Bandit achieves 71.4\% detection across 11 pattern-based CWEs but 0\% across 19 logic-flaw CWEs. LLM-generated tests provide their primary value on logic-flaw CWEs where static analysis is blind---authentication bypass (CWE-306), CSRF (CWE-352), IDOR (CWE-639), and SSRF (CWE-918) can only be detected through behavioral testing.}

\fix{
Cross-model per-CWE analysis reveals consistent patterns. \textbf{CWE-89 (SQL Injection)} achieves 100\% SMS across all models with qualifying samples---a universally well-detected vulnerability. \textbf{CWE-732 (Weak Permissions)} is consistently high (71--100\% SMS) for most models except GLM-5 (0\%). \textbf{CWE-611 (XXE)} is a universal blind spot: 0\% SMS across all eight models, as tests consistently crash due to API mismatches between \texttt{defusedxml} and standard XML parsers. \textbf{CWE-95 (Eval Injection)} is also near-zero across all models. \textbf{CWE-918 (SSRF)} shows 0\% SMS for most models, with only Kimi-K2.5 (21.4\%) and Qwen2.5-Coder (28.6\%) achieving any detection. Model coverage breadth varies: GLM-5 covers 27 CWEs (broadest presence) but many at 0\% SMS, while Kimi-K2.5 covers 25 CWEs with generally higher SMS values.
}

%% FIGURE: Per-CWE SMS Heatmap
\begin{figure}[t]
  \centering
  \fbox{\parbox{0.93\columnwidth}{\centering\vspace{1.2cm}\textbf{Figure: Per-CWE SMS Heatmap}\\\smallskip
  \small Heatmap with CWE categories on y-axis (grouped: Injection, Crypto, Auth/Access, Other) and models + Bandit on x-axis. Cell color intensity = SMS value (darker = higher). Annotate cells where Bandit $>$ LLM SMS and vice versa. Alternative: radar/spider chart with one polygon per model overlaid, axes = CWE categories, to show model-specific strengths and blind spots.\vspace{1.2cm}}}
  \caption{Per-CWE Security Mutation Score across models and Bandit baseline. \fix{Pattern-based CWEs (top) vs.\ logic-flaw CWEs (bottom) reveal distinct detection profiles.}}
  \label{fig:cwe_heatmap}
\end{figure}

\smallskip\noindent\fbox{\parbox{0.95\columnwidth}{\fix{\textbf{Finding 3:} Static analysis (Bandit) achieves 71.4\% detection on 11 pattern-based CWEs but 0\% on 19 logic-flaw CWEs. LLM-generated tests provide complementary coverage on logic-flaw vulnerabilities (authentication, authorization, CSRF, SSRF) where static analysis is structurally blind.}}}\smallskip

\fix{
\subsection{RQ4: Prompt Sensitivity and Baselines}

\textbf{Prompt ablation.} Table~\ref{tab:ablation} compares SMS across three prompt variants (no hint, CWE~ID only, full context) to quantify how much security test quality depends on explicit vulnerability guidance. All results use OpenAI-judged classification.

\begin{table}[t]
\caption{Prompt ablation: SMS (\%) by prompt variant.}
\label{tab:ablation}
\small
\begin{tabular}{lrrr}
\toprule
\textbf{Model} & \textbf{No hint} & \textbf{CWE ID} & \textbf{Full context} \\
\midrule
GPT-oss-120B & 0.0 & 52.3 & 43.2 \\
GPT-5.2 & 31.2 & 54.5 & 38.3 \\
GPT-5-mini & 20.7 & 38.1 & 39.8 \\
Kimi-K2.5 & 30.1 & 57.5 & 45.7 \\
Qwen2.5-Coder 14B & 12.9 & 14.9 & 58.6 \\
Qwen3-Coder 30B & 35.5 & 42.5 & 42.3 \\
DeepSeek-Coder-V2 & 7.1 & 7.5 & 10.9 \\
GLM-5 & 25.6 & 34.3 & 27.4 \\
\bottomrule
\end{tabular}
\end{table}

The relationship between prompt detail and SMS is non-monotonic. Surprisingly, \emph{CWE-ID-only} prompts yield the highest SMS for 4 of 8 models (GPT-oss: 52.3\%, GPT-5.2: 54.5\%, Kimi: 57.5\%, Qwen3: 42.5\%), outperforming full-context prompts. This suggests that concise CWE identification may focus models on the vulnerability class without overwhelming them with implementation details. Conversely, Qwen2.5-Coder shows the opposite pattern: full context (58.6\%) dramatically outperforms CWE-ID (14.9\%), suggesting this model benefits from detailed guidance.

No-hint prompts consistently degrade SMS, with GPT-oss-120B dropping to 0\% SMS (producing no security-aware assertions without any vulnerability hint) and DeepSeek-Coder-V2 achieving only 7.1\%. The average SMS delta between no-hint and best-variant is 24.8 percentage points, confirming that explicit vulnerability guidance is essential for security-aware test generation.

\textbf{Reference test baseline.} Expert-written reference tests achieve MS of 57.7\% and SMS of 47.6\% with 889 semantic kills out of 1,869 mutants and a perfect 100\% secure-pass rate. The best LLM (Qwen2.5-Coder at 58.6\% SMS) exceeds reference test SMS, but on only 16.8\% of samples. When accounting for secure-pass rate, the \emph{effective SMS} (SMS $\times$ SPR) tells a different story: reference tests achieve 47.6\% effective SMS, while the best LLM effective SMS is Qwen3-Coder at $42.3\% \times 46.6\% = 19.7\%$.
}

\smallskip\noindent\fbox{\parbox{0.95\columnwidth}{\fix{\textbf{Finding 4:} CWE-ID-only prompts outperform full-context prompts on SMS for half the models, while no-hint prompts severely degrade performance (0--35.5\% SMS). Expert reference tests achieve 47.6\% EffSMS; the best LLM reaches only 19.7\% EffSMS when accounting for test validity.}}}\smallskip


%% ====================================================================
%% VI. DISCUSSION
%% ====================================================================
\section{Discussion}\label{sec:discussion}

\subsection{Implications for LLM Evaluation}

\fix{Our findings carry significant implications for how the community evaluates LLM security capabilities. The 2.4$\times$ gap between expert-written reference tests (47.6\% EffSMS) and the best LLM (19.7\% EffSMS) demonstrates that LLMs remain far from expert-level security test generation. This gap compounds from two sources: most generated tests fail on secure code (SPR of 7--47\%), and among valid tests, a substantial proportion of kills reflect functional or crash behavior rather than security awareness (MS/SMS inflation of 2.2$\times$). A model reporting 90\% MS may deliver only 5\% EffSMS in practice (GPT-5.2: 88.7\% MS $\rightarrow$ 5.1\% EffSMS). Benchmarks that report only pass/fail or mutation scores without accounting for test validity and kill classification risk overestimating model capabilities by an order of magnitude.}

The \fix{dominance of functional kills} suggests that LLM-generated tests often detect \fix{behavioral changes caused by mutations through correctness assertions} rather than security vulnerabilities. This parallels findings in traditional mutation testing, where Mouelhi et al.~\cite{mouelhi2007mutation} demonstrated that functional tests achieving high code coverage fail to achieve high security mutation scores. Our results extend this insight to the LLM context: models that are proficient at generating functional tests may nevertheless lack genuine security reasoning capability.


\subsection{Static Analysis Comparison}

The Bandit~\cite{bandit2024} baseline provides important context. On pattern-based vulnerabilities with clear syntactic signatures (e.g., \texttt{shell=True}, string-formatted SQL), static analysis achieves high detection rates (\fix{71.4\%} across static-analyzable CWEs). However, \fix{56\%} of our dataset consists of logic-flaw CWEs where Bandit achieves \fix{0\%} detection. This demonstrates that mutation testing provides essential coverage that static analysis cannot: it evaluates whether tests detect \emph{behavioral} vulnerabilities, not just \emph{syntactic} anti-patterns.

%% FIGURE: Static vs Dynamic Detection Comparison
\begin{figure}[t]
  \centering
  \fbox{\parbox{0.93\columnwidth}{\centering\vspace{1.2cm}\textbf{Figure: Static vs.\ Dynamic Detection Coverage}\\\smallskip
  \small Grouped bar chart with CWEs on x-axis grouped into ``Pattern-based'' \fix{(11 CWEs)} and ``Logic-flaw'' \fix{(19 CWEs)}. Two bars per CWE: Bandit detection rate (blue) and best-LLM SMS (orange). Shows Bandit dominating pattern-based CWEs while LLMs provide coverage on logic-flaw CWEs. Add horizontal dashed lines for Bandit average (\fix{22.7\%}) and overall LLM SMS average.\vspace{1.2cm}}}
  \caption{Complementary coverage: Bandit excels on pattern-based vulnerabilities while LLM-generated tests provide value on logic-flaw CWEs where static analysis achieves \fix{0\%} detection.}
  \label{fig:static_vs_dynamic}
\end{figure}

\fix{On pattern-based CWEs where Bandit achieves 71.4\% detection (CWE-89, 22, 94, 327, 798, and 6 others), the best LLMs match or exceed Bandit on several: CWE-89 (all models achieve 100\% SMS vs.\ Bandit 71.4\%), CWE-327 (GPT-5.2 achieves 94.1\% SMS), and CWE-798 (GPT-5-mini achieves 100\% SMS). However, on CWE-22 (Path Traversal), most LLMs achieve 0\% SMS while Bandit detects 71.4\%. On the 19 logic-flaw CWEs where Bandit achieves 0\%, LLMs provide unique value: CWE-306 (Kimi-K2.5 achieves 75.0\% SMS), CWE-352 (DeepSeek-V2 achieves 66.7\%), CWE-732 (GPT-5-mini achieves 97.2\%), and CWE-915 (Kimi-K2.5 achieves 73.3\%). This confirms complementary coverage: the optimal strategy combines static analysis for syntactic patterns with LLM-generated tests for behavioral vulnerabilities.}

\fix{
\textbf{Static analysis on mutants vs.\ originals.} A natural question is whether running static analysis on \emph{mutated} code rather than insecure originals would change these findings. For syntactic operators (PSQLI, WEAKCRYPTO, NOCERTVALID), the result is equivalent---the mutant introduces the same syntactic pattern Bandit detects in the insecure original (e.g., string-formatted SQL, \texttt{verify=False}). For behavioral operators (SSRF, MISSINGAUTH, IDOR, CSRF\_REMOVE), Bandit achieves 0\% detection on both originals and mutants by construction: these vulnerabilities manifest as \emph{missing logic}, not as detectable syntactic patterns. Running Bandit on mutants would therefore not change our complementarity finding---static analysis covers syntactic CWEs regardless of whether the input is an original or a mutant, and is structurally blind to logic-flaw CWEs in either case.
}

\subsection{Why Current Metrics Overstate Capability}

\fix{The inflation mechanism operates through three channels. First, \textbf{functional kills} arise when mutations alter control-flow behavior caught by behavioral-contract assertions---for example, a \textsc{MISSINGAUTH} mutation bypassing an access check so that an expected \texttt{PermissionError} is no longer raised. The test detects the mutation through an exception-expectation assertion, not a security assertion. Second, \textbf{incidental kills} arise when mutations change data values that a non-security equality assertion happens to check---for example, a \textsc{WEAKCRYPTO} mutation changing hash output length from 64 to 32 characters, caught by \texttt{assert len(result) == 64} rather than a cryptographic-strength check. Third, \textbf{crash kills} arise when mutations break syntactic or type-level properties, causing runtime exceptions (ImportError, TypeError, NameError) before any assertion is evaluated. None of these three channels require security awareness.}

\fix{
\textbf{Functional kill example.} For CWE-95 (Eval Injection), a GPT-5-mini test expects \texttt{ValueError} when calling a formula evaluation function. The EVALINJECT mutant replaces \texttt{ast.parse} with \texttt{eval()}, which silently evaluates the input instead of raising an error. The test assertion \texttt{DID NOT RAISE <class 'ValueError'>} catches the mutation through expected-exception checking, not through any assertion about code injection. This functional kill inflates MS without contributing to SMS.

\textbf{Crash kill example.} For CWE-327 (Weak Cryptography), a Qwen3-Coder test calls a hashing function. The WEAKCRYPTO mutant removes the hash computation line entirely, causing \texttt{NameError: name 'v\_c9449f\_attachment' is not defined}. The test never reaches any assertion---the mutant is killed by structural code breakage, not by security-aware detection.
}

\textcolor{red}{
\subsection{SMS Robustness and Gaming Risks}



A potential concern is that SMS, being based on keyword matching in assertion messages, could be gamed by LLMs that learn to insert security-relevant terms without genuine understanding (e.g., \texttt{assert ``sql injection'' in result} as a vacuous assertion). We identify three structural mitigations already present in our design: (1)~operator-aware keywords prevent cross-CWE gaming---inserting ``sql injection'' only boosts SMS for \textsc{PSQLI} mutants, not for \textsc{WEAKCRYPTO} or \textsc{DESERIAL} mutants; (2)~the kill predicate requires tests to \emph{pass} on secure code and \emph{fail} on the mutant, so keyword insertion alone is insufficient without behavioral detection; and (3)~our mock objects track security-relevant execution state (e.g., \texttt{MockDatabase.last\_query}, \texttt{MockSubprocess.last\_shell}, \texttt{MockFileSystem.last\_path}) that could support \emph{observability-based} classification in future work.



An observability-based extension would classify a kill as semantic not based on assertion \emph{text} but on whether the test inspects mock state attributes that reflect the security property under test. For example, a test that asserts \texttt{subprocess.last\_shell == False} after calling the function directly checks whether shell injection is prevented---regardless of what text appears in the assertion message. We leave this strengthening to future work, noting that the infrastructure (mock state tracking) is already in place.



\subsection{Mock Environment Validity}



\fix{Heavy reliance on mocks risks ecological validity: test behavior in mocked environments may diverge from production execution. Our reference tests provide a partial validity check: expert-written tests achieve 57.7\% MS and 47.6\% SMS in the mocked environment, with 100\% secure-pass rate, demonstrating that the mock infrastructure supports meaningful security testing. The 82.5\% semantic kill proportion for reference tests (889/1,078 kills) confirms that the mock environment enables genuine security-aware assertions rather than systematically biasing toward crash or functional kills. The primary sources of mock-induced divergence are (1)~import interception via \texttt{sys.modules} that may prevent tests from exercising real library behavior---notably, CWE-611 (XXE) achieves 0\% SMS across all models partly due to API differences between mocked and real XML parsers---and (2)~mock return values that may not capture all edge cases of real APIs. We note that mocks are a necessary safety layer---without them, LLM-generated tests could execute arbitrary system commands---and that the trade-off between safety and ecological validity is inherent to any sandboxed evaluation.}

}

\subsection{Recommendations}

Based on our findings, we recommend that: (1) evaluations of LLM security test generation should decompose kills by type rather than reporting aggregate mutation scores; (2) benchmark datasets should include pre-generated mutants for reproducibility; and (3) security test quality metrics should weight semantic kills substantially higher than crash or incidental kills. For practitioners using LLMs to generate security tests, our results suggest generated tests should be reviewed for security-specific assertions rather than trusted based on pass/fail behavior alone.


\section{\color{red}Need to be modified in the method and results}\label{sec:threats}

%% ====================================================================
%% VIII. THREATS TO VALIDITY
%% ====================================================================
\section{Threats to Validity}\label{sec:threats2}

\textbf{Internal validity.} Kill classification relies on error message pattern matching, which may misclassify kills where security-relevant terms appear incidentally or where genuine security assertions use non-standard terminology (e.g., checking mock call histories or return codes without keyword-bearing messages). We mitigate this with operator-specific keyword lists rather than generic patterns, \fix{and validate accuracy using an LLM-as-Judge (GPT-5.4) to reclassify heuristic labels}. \fix{Tests that assert security properties via mock state inspection (e.g., \texttt{db.last\_params is not None}) without explicit keyword usage may be classified as functional rather than semantic by the heuristic---a known limitation that underestimates SMS, mitigated by the LLM-as-Judge reclassification.} Mutation operator coverage may not capture all realistic vulnerability patterns\fix{; expanding operator coverage to additional CWE subcategories remains future work}.

\textbf{External validity.} SecMutBench covers \fix{30} CWEs in Python, limiting generalization to other vulnerability types and languages. Our CWE selection (Section~\ref{sec:framework}) prioritizes depth over breadth, ensuring each CWE has verified operators and classification keywords, but the dataset may not represent the full diversity of real-world vulnerability patterns. The \fix{5.51}~mutants per program reflects our operator design targeting specific vulnerability patterns rather than exhaustive syntactic variation; while adequate for measuring security-aware detection, it means that survival-based analyses (equivalent mutant estimation, operator subsumption) require caution. Our contamination prevention pipeline mitigates training-data overlap through structural perturbation, identifier renaming, temporal filtering, and n-gram auditing, but cannot guarantee zero contamination for all models---particularly those trained on datasets not available for audit.

\textbf{Construct validity.} SMS assumes that semantic kills indicate genuine security awareness, but a test could contain security keywords without meaningful understanding. We discuss this gaming risk and structural mitigations in Section~\ref{sec:discussion}. The mock-based execution environment may affect test behavior differently than production environments\fix{; we acknowledge this trade-off between safety and ecological validity (Section~\ref{sec:discussion})}. Equivalent mutants, while structurally unlikely for security operators that alter behavioral properties (Section~3.2), may occur for operators targeting validation removal on unreachable paths; we report the executable mutant ratio to bound this effect.

\fix{
\textbf{Mock ecological validity.}
The sandboxed execution environment introduces a known validity gap
for operators where the mock API diverges from the real library.
The clearest case is CWE-611 (XXE): the \textsc{XXE} operator
replaces \texttt{defusedxml} with standard \texttt{xml.etree},
but MockXMLParser exposes a different API surface than either
library. Across all eight models, \textsc{XXE} achieves 0\% SMS
and a high crash rate (Table~\ref{tab:ira}),
confirming that tests fail before reaching any assertion rather
than detecting the vulnerability. We flag CWE-611 results as
unreliable and exclude them from per-CWE ranking in
Section~\ref{sec:results}. More broadly, the mock environment is a
necessary safety constraint---without it, LLM-generated tests could
execute arbitrary system calls---but practitioners should treat
mock-dependent SMS values as conservative lower bounds relative to
production execution. Quantifying the full ecological validity gap
(e.g., by comparing mocked vs.\ real-but-sandboxed execution)
remains future work.
}



%% ====================================================================
%% IX. CONCLUSION
%% ====================================================================
\section{Conclusion}\label{sec:conclusion}

We presented SecMutBench, a mutation-based benchmark for evaluating LLM-generated security tests. Through \fix{25} security mutation operators applied to \fix{339} Python programs across \fix{30} CWE categories\fix{, producing 1,869 mutants}, we demonstrated that traditional mutation scores substantially overstate LLM security testing capability. The Security Mutation Score (SMS) metric and operator-aware kill classification reveal that \fix{functional assertions and runtime crashes, rather than genuine security awareness, account for a substantial proportion of mutant kills}.

\fix{Our evaluation of eight LLMs reveals that the best model achieves only 19.7\% EffSMS (Qwen3-Coder~30B) compared to 47.6\% for expert-written reference tests---a 2.4$\times$ gap. This gap compounds from two sources: only 7--47\% of generated tests pass on secure code, and among valid tests, traditional MS overstates security capability by 2.2$\times$ on average. A model reporting 88.7\% MS (GPT-5.2) delivers only 5.1\% EffSMS in practice.}

Comparison with Bandit shows that static analysis and mutation testing provide complementary coverage, with static tools excelling on pattern-based vulnerabilities \fix{(71.4\% detection on 11 CWEs)} while mutation testing captures logic-flaw detection capability \fix{(19 CWEs where Bandit achieves 0\%)}.

\fix{SecMutBench is publicly available with pre-generated mutants, executable evaluation framework, 11 CWE-specific mock environments, and a sandboxed execution layer to support reproducible security test evaluation research. The code, evaluation pipeline, and replication package are available at \url{https://github.com/Mars-2030/secmutbench}; the dataset (339 samples, 1,869 mutants, 30 CWEs) is hosted at \url{https://huggingface.co/datasets/Mars203020/secmutbench}.}

\textcolor{red}{
\textbf{Future work.} We plan to expand CWE coverage, add support for additional languages, and investigate few-shot and chain-of-thought prompting effects on SMS. We also plan to explore \emph{observability-based kill classification} that leverages mock state tracking (e.g., \texttt{MockDatabase.last\_params}, \texttt{MockSubprocess.last\_shell}) rather than assertion message keywords, reducing susceptibility to superficial keyword insertion. Following MuTAP~\cite{dakhel2023mutap} and MUTGEN~\cite{mutgen2025}, we will investigate iterative refinement where surviving security mutants are fed back to models as prompts, measuring whether SMS improves with mutation-guided feedback. Finally, we plan to evaluate coding agents (e.g., Cursor, GitHub Copilot) in realistic workflows to assess the practical impact of our benchmark findings.
}


%%
%% The acknowledgments section is defined using the "acks" environment
%% (and NOT an unnumbered section). This ensures the proper
%% identification of the section in the article metadata, and the
%% consistent spelling of the heading.
\begin{acks}
To Robert, for the bagels and explaining CMYK and color spaces.
\end{acks}

%%
%% The next two lines define the bibliography style to be used, and
%% the bibliography file.
\bibliographystyle{ACM-Reference-Format}
\bibliography{ref}


%%
%% If your work has an appendix, this is the place to put it.
\appendix




\end{document}
\endinput
%%
%% End of file `sample-sigconf.tex'.
